%%%%%%%%%%%%%%%%%%%%%%%%%%%%% Define Article %%%%%%%%%%%%%%%%%%%%%%%%%%%%%%%%%%
\documentclass{article}
%%%%%%%%%%%%%%%%%%%%%%%%%%%%%%%%%%%%%%%%%%%%%%%%%%%%%%%%%%%%%%%%%%%%%%%%%%%%%%%

%%%%%%%%%%%%%%%%%%%%%%%%%%%%% Using Packages %%%%%%%%%%%%%%%%%%%%%%%%%%%%%%%%%%
\usepackage{geometry}
\usepackage{graphicx}
\usepackage{amssymb}
\usepackage{amsmath}
\usepackage{amsthm}
\usepackage{empheq}
\usepackage{mdframed}
\usepackage{booktabs}
\usepackage{lipsum}
\usepackage{graphicx}
\usepackage{color}
\usepackage{psfrag}
\usepackage{pgfplots}
\usepackage{bm}
%%%%%%%%%%%%%%%%%%%%%%%%%%%%%%%%%%%%%%%%%%%%%%%%%%%%%%%%%%%%%%%%%%%%%%%%%%%%%%%

% Other Settings

%%%%%%%%%%%%%%%%%%%%%%%%%% Page Setting %%%%%%%%%%%%%%%%%%%%%%%%%%%%%%%%%%%%%%%
\geometry{a4paper}

%%%%%%%%%%%%%%%%%%%%%%%%%% Define some useful colors %%%%%%%%%%%%%%%%%%%%%%%%%%
\definecolor{ocre}{RGB}{243,102,25}
\definecolor{mygray}{RGB}{243,243,244}
\definecolor{deepGreen}{RGB}{26,111,0}
\definecolor{shallowGreen}{RGB}{235,255,255}
\definecolor{deepBlue}{RGB}{61,124,222}
\definecolor{shallowBlue}{RGB}{235,249,255}
%%%%%%%%%%%%%%%%%%%%%%%%%%%%%%%%%%%%%%%%%%%%%%%%%%%%%%%%%%%%%%%%%%%%%%%%%%%%%%%

%%%%%%%%%%%%%%%%%%%%%%%%%% Define an orangebox command %%%%%%%%%%%%%%%%%%%%%%%%
\newcommand\orangebox[1]{\fcolorbox{ocre}{mygray}{\hspace{1em}#1\hspace{1em}}}
%%%%%%%%%%%%%%%%%%%%%%%%%%%%%%%%%%%%%%%%%%%%%%%%%%%%%%%%%%%%%%%%%%%%%%%%%%%%%%%

%%%%%%%%%%%%%%%%%%%%%%%%%%%% English Environments %%%%%%%%%%%%%%%%%%%%%%%%%%%%%
\newtheoremstyle{mytheoremstyle}{3pt}{3pt}{\normalfont}{0cm}{\rmfamily\bfseries}{}{1em}{{\color{black}\thmname{#1}~\thmnumber{#2}}\thmnote{\,--\,#3}}
\newtheoremstyle{myproblemstyle}{3pt}{3pt}{\normalfont}{0cm}{\rmfamily\bfseries}{}{1em}{{\color{black}\thmname{#1}~\thmnumber{#2}}\thmnote{\,--\,#3}}
\theoremstyle{mytheoremstyle}
\newmdtheoremenv[linewidth=1pt,backgroundcolor=shallowGreen,linecolor=deepGreen,leftmargin=0pt,innerleftmargin=20pt,innerrightmargin=20pt,]{theorem}{Theorem}[section]
\theoremstyle{mytheoremstyle}
\newmdtheoremenv[linewidth=1pt,backgroundcolor=shallowBlue,linecolor=deepBlue,leftmargin=0pt,innerleftmargin=20pt,innerrightmargin=20pt,]{definition}{Definition}[section]
\theoremstyle{myproblemstyle}
\newmdtheoremenv[linecolor=black,leftmargin=0pt,innerleftmargin=10pt,innerrightmargin=10pt,]{problem}{Problem}[section]
%%%%%%%%%%%%%%%%%%%%%%%%%%%%%%%%%%%%%%%%%%%%%%%%%%%%%%%%%%%%%%%%%%%%%%%%%%%%%%%

%%%%%%%%%%%%%%%%%%%%%%%%%%%%%%% Plotting Settings %%%%%%%%%%%%%%%%%%%%%%%%%%%%%
\usepgfplotslibrary{colorbrewer}
\pgfplotsset{width=8cm,compat=1.9}
%%%%%%%%%%%%%%%%%%%%%%%%%%%%%%%%%%%%%%%%%%%%%%%%%%%%%%%%%%%%%%%%%%%%%%%%%%%%%%%

%%%%%%%%%%%%%%%%%%%%%%%%%%%%%%% Title & Author %%%%%%%%%%%%%%%%%%%%%%%%%%%%%%%%
\title{Chinese Remainder Theorem}
\author{Alston Yam}
\date{}
\parskip=5pt
\parindent=0pt
%%%%%%%%%%%%%%%%%%%%%%%%%%%%%%%%%%%%%%%%%%%%%%%%%%%%%%%%%%%%%%%%%%%%%%%%%%%%%%%

\begin{document}
    \maketitle

    \section{Introduction}
    Lets have a look at values of $x$ that satisfy the following equations:
    \begin{align*}
        x &\equiv 2 \pmod{7} \\
        x &\equiv 1 \pmod{3}
    \end{align*}

    We can think about solutions satisfying each of the individual congruences. To the first one: 
    \[A = \{\ldots, -19, -12, -5, 2, 9, 16, 23, \ldots\}\]

    And the second one:
    \[B = \{\ldots, -8, -5, -2, 1, 4, 7, 10, \ldots\}\]

    So we simply see that all the possible values of $x$ are $x \in A \cup B$.

    \subsection{Other possible cases}
    There are also some other cases to think about: 
    \begin{align*}
        x &\equiv 1 \pmod{3} \\
        x &\equiv 1 \pmod{6}
    \end{align*}

    Or consider this one:
    \begin{align*}
        x &\equiv 1 \pmod{6} \\
        x &\equiv 2 \pmod{9}
    \end{align*}

    What happens in each case when we try to find the general solution? These experiments are left as an exercise to the reader. 
    \pagebreak


    \section{The theorem itself}
    \begin{theorem}[Chinese Remainder Theorem]
        Let $a_1, a_2, \ldots, a_n$ be integers, and $b_1, b_2, \ldots, b_n$ be pairwise coprime integers. Then the system of equations
        \begin{align*}
            x &\equiv a_1 \pmod{b_1} \\
            x &\equiv a_2 \pmod{b_2} \\
            \vdots \\
            x &\equiv a_n \pmod{b_n}
        \end{align*}
        Has a $\textit{unique}$ solution under $\pmod{b_1b_2 \ldots b_n}$.
    \end{theorem}

    \begin{proof}[Proof of CRT]

        We split the problem into uniqueness and existence:

        \textbf{Uniqueness}: Assume that there are two solutions $x, y$ that satisfy all the congruences. Then we must get that $x-y$ is divisible by $b_i$ $\forall i$. However we have the condition that $b_i$ are pairwise coprime, so $x-y$ must be divisible by $N = b_1b_2 \ldots b_n$, and as a result, $x \equiv y \pmod{N}$, contradiction.

        \textbf{Existence}: We will proceed by induction on n. 

        The case $n=1$ is trivial, so we start from $n=2$. Consider the congruences 
        \begin{align*}
            x &\equiv a_1 \pmod{b_1} \\
            x &\equiv a_2 \pmod{b_2}
        \end{align*}

        By Bezout's identity, we have there exist integers $m_1, m_2$ such that $m_1b_1 + m_2b_2 = 1$. Lets construct $x = a_2m_1b_1 + a_1m_2b_2$. We will see that this $x$ is the desired x: 
        \[x = a_2(1-m_2b_2) + a_1m_2b_2 = a_2 + m_2b_2(a_1 - a_2) \equiv a_2 \pmod{b_2}\]

        Then the other congruence is satisfied similarly. 

        Assume that there exists a unique $x$ such that $n$ of these congruences hold true simultaneously. 

        Consider $n+1$ congruences:
        \begin{align*}
            x &\equiv a_1 \pmod{b_1} \\
            x &\equiv a_2 \pmod{b_2} \\ 
            \vdots \\
            x &\equiv a_n \pmod{b_n} \\ 
            x &\equiv a_{n+1} \pmod{b_{n+1}}
        \end{align*}
        We can combine these like so: 

        \begin{align*}
            x &\equiv a_{1, 2, \ldots n} \pmod{b_{1, 2, \ldots n}} \\
            x &\equiv a_{n+1} \pmod{b_{n+1}}
        \end{align*}

        And here we have our $n=2$ base case which we know to work.




        
    \end{proof}

    \section{Applications of the theorem}
    
    This theorem is extremely useful when we want to construct an $x$ that satisfies many properties. 

    \begin{problem}
        Show that for $c \in \mathbb{Z}$ and a prime $p$, the congruence $x^x \equiv c \pmod{p}$ has a solution. 
    \end{problem}

    The idea is that we are looking at a prime modulus, so FLT probably is going to be useful. Note that $x^{p-1} \equiv 1 \pmod{p}$, so the modulus cycles every $p-1$. 

    Consider the system of congruences: 

    \begin{align*}
        x &\equiv 1 \pmod{p-1} \\
        x &\equiv c \pmod{p}
    \end{align*}

    Then we know that since $\gcd{(p, p-1)} = 1$, there is a solution to $x$ by CRT, and hence we have $x^x \equiv x^{x \pmod{p-1}} \equiv x^1 \equiv c \pmod{p}$ and so we are done.

    \begin{problem}
        What are the last 2 digits of $49^{19}$?
    \end{problem}

    So we are interested in $49^{19} \pmod{100}$. Notice that $100 = 25 \times 4$ and $\gcd(25, 4) = 1$. This means that we only really need to find the $x$ such that 
    \begin{align*}
        x &\equiv 49^{19} \pmod{25} \\
        x &\equiv 49^{19} \pmod{4}
    \end{align*}

    We can compute each of the congruences and turn it into 
    \begin{align*}
        x &\equiv -1 \pmod{25} \\
        x &\equiv 1 \pmod{4}
    \end{align*}

    Which, by CRT, has unique solution $x \equiv ((-1)(4)(19) + (1)(25)(1)) \equiv 49 \pmod{100}$, so we know that the last 2 digits of $x$ are 49. 









    
\end{document}